% !TeX program = pdflatex
% ALIUS Bulletin standalone facsimile source.
% Compile from the ALIUS-bulletin project root for exact visual fidelity.
% If the shared PDF asset is absent, this file falls back to extracted text.
\providecommand{\ALIUSRootPrefix}{}

\ifdefined\ALIUSIssueBuild
  \IfFileExists{\ALIUSRootPrefix Shared-assets/original-pdfs/Issue04/alius-bulletin-issue-04-tribute-to-martin-fortier.pdf}{%
    \includepdf[pages=-,fitpaper=true]{\ALIUSRootPrefix Shared-assets/original-pdfs/Issue04/alius-bulletin-issue-04-tribute-to-martin-fortier.pdf}%
  }{%
    \typeout{ALIUS warning: missing source PDF Shared-assets/original-pdfs/Issue04/alius-bulletin-issue-04-tribute-to-martin-fortier.pdf}%
  }%
  \expandafter\endinput
\fi

\documentclass[a4paper,11pt]{article}
\usepackage[margin=22mm]{geometry}
\usepackage{pdfpages}
\usepackage[hidelinks]{hyperref}
\usepackage{microtype}
\usepackage{parskip}
\title{A tribute to Martin E. Fortier}
\author{by ALIUS Research group}
\date{ALIUS Bulletin 4}
\begin{document}
\IfFileExists{\ALIUSRootPrefix Shared-assets/original-pdfs/Issue04/alius-bulletin-issue-04-tribute-to-martin-fortier.pdf}{%
  \includepdf[pages=-,fitpaper=true]{\ALIUSRootPrefix Shared-assets/original-pdfs/Issue04/alius-bulletin-issue-04-tribute-to-martin-fortier.pdf}%
}{%
  \maketitle
  \section*{Metadata}
  \textbf{Piece type:} tribute\par
  \textbf{DOI:} 10.34700/aqmn-r841\par
  \textbf{Keywords:} hallucinogens, bayesian constructivism, neuropharmacology, psychedelics, religion, consciousness, homo neurochemicus, HUTHAC\par
  \section*{Abstract}
  Martin Fortier has been a wonderful friend and a brilliant colleague, gone far too soon to realize his extraordinary potential despite his many precocious achievements. More than anyone, he embodied the spirit of interdisciplinarity. His passion for the scientific study of consciousness was inspiring, altering the trajectory of many researchers in our network and beyond. This text is intended to provide a faithful, up-to-date, and accurate representation of Martin's wide-ranging research by adapting in large part from Martin's own summary of his work written about a year before his passing.\par
  \section*{Extracted Text}
  \begingroup
  \small
  \noindent ALIUS Bulletin n4 (2020) aliusresearch.org/bulletin Tribute to Martin Fortier Abstract Martin Fortier has been a wonderful friend and a brilliant colleague, gone far too soon to realize his extraordinary potential despite his many precocious achievements. More than anyone, he embodied the spirit of interdisciplinarity. His passion for the scientific study of consciousness was inspiring, altering the trajectory of many researchers in our network and beyond. This text is intended to provide a faithful, up-to-date, and accurate representation of Martin's wide- ranging research by adapting in large part from Martin's own summary of his work written about a year before his passing. keywords: hallucinogens, bayesian constructivism, neuropharmacology, psychedelics, religion, consciousness, homo neurochemicus, HUTHAC On April 11th, 2020, Martin Fortier tragically passed away after a long and harrowing battle with cancer. He was thirty years old. We mourn the loss of a wonderful friend and a brilliant colleague, gone far too soon to realize his extraordinary potential despite his many precocious achievements. Martin co-founded ALIUS in 2016, and has been a pillar of our group since its inception. More than anyone, he embodied the spirit of interdisciplinarity at the core of ALIUS' mission. His unique academic trajectory reveals in that respect: after a Bachelor's degree in philosophy at the Sorbonne, Martin obtained two Master's degrees in parallel - one in philosophy, and another in anthropology. He then started a Ph.D. in cognitive science at the Institut Jean Nicod, under the supervision of Jerome Dokic. His dissertation, which remains unfinished, is an impressive tour de force seamlessly connecting ideas and data from philosophy, anthropology, psychology, neuroscience, and pharmacology. Martin's ideal of interdisciplinarity was rooted in his aversion for empirical blindspots, partial evidence, and culturally-grounded intuitions. He encouraged philosophers to seek the support of empirical data and refrain Cite as: ALIUS Research Group (2020). Tribute to Martin Fortier. ALIUS Bulletin, 4, 1-14, https://doi.org/10.34700/aqmn-r841 ALIUS Research Group aliusresearchgroup@gmail.com ALIUS Research Group - Tribute to Martin Fortier ALIUS Bulletin n4 (2020) aliusresearch.org/bulletin 2 from making underspecified or speculative claims about topics that can be empirically determined. He invited neuroscientists and psychologists to be more aware of the narrow cultural background of their experiments, that often recruit "WEIRD" participants - those raised in a society that is Western, Educated, Industrialized, Rich, and Democratic. Finally, he incited anthropologists to relate their theories to empirically-grounded hypotheses about the cognitive architecture of the human mind. Martin strived to apply these injunctions to his own research. His life and work embodied many intellectual virtues rarely conjoined in a single individual to such a high degree: curiosity, passion, meticulousness, rigor, tenacity. Those who knew him remember how he could rekindle one's interest in research or spur ideas for a whole new project within a single conversation - so infectious was his love of learning and his enthusiasm for charting the human mind. When he set out to study a topic, he would get consumed by reading about it and leave no stone unturned. In this quest for encyclopedic knowledge, he devoured articles and books from many disciplines, and had an impressive capacity for elegantly synthesizing them in order to yield new insights. His research combined a number of methods and tools borrowed from his interdisciplinary expertise: conceptual analysis from analytic philosophy; theorization based on existing anthropological, psychological, neuroscientific and biological data; collection of new experimental data, and collection of ethnographic data. His fieldwork career included work with Shipibo-Konibo communities of the Ucayali region in the Peruvian Amazon, as well as Huni Kuin (Cashinahua) communities on the Peruvian side of the Purus region. Martin liked to organize his many research interests within two broad categories. In the first category was his research on the structure of abstract knowledge in indigenous Amazonian cultures; in the second, his research on hallucinogens, culture, and neurobiology. One of the key ideas behind Martin's work on abstract knowledge in Amazonian cultures is Bayesian constructivism. There is a long-lasting controversy in psychology between nativists, who think abstract knowledge is innately present in humans, and empiricists, who believe it is simply ALIUS Research Group - Tribute to Martin Fortier ALIUS Bulletin n4 (2020) aliusresearch.org/bulletin 3 abstracted from regularities detected by the senses, or socially constructed. Martin favored a third proposal - Bayesian constructivism - according to which abstract knowledge is more than the simple accumulation of sensory data, but is nonetheless gradually constructed through development rather than being innately given or directly taught. One implication of this view is that different abstract structures may be constructed given different kinds of experience. Martin's work explored the fruitfulness of the Bayesian constructivist framework in studying how Amazonian indigenous people categorize and reason about the world. One of Martin's targets was psychological essentialism, the view that external properties are explained by inner ones, and that identity is defined by inner and fixed properties rather than superficial and malleable ones. Many authors claim that essentialism is found in all cultures across the world. Martin's work challenged this claim. In particular, he argued that Cashinahua people use non-essentialist forms of reasoning. For them, many inner properties are explained by external ones and are thus highly malleable. He claimed that Cashinahua non-essentialism obtains both in the domains of ethnozoology (reasoning about different species of animals) and ethnosociology (reasoning about different ethnic groups). Martin further considered the implications of the existence of non- essentialist reasoning in Amazonian cultures for anthropology. Anthropologists such as Descola and Viveiros de Castro have proposed that what is special in the way Amazonian people think is that they ascribe mental and bodily properties differently than we do (for example, as compared to us, they "over-ascribe" mental properties to animals and plants). Martin's work on non-essentialism suggests instead that what differentiates Amazonian thinking from our own thinking is not so much how mental or bodily properties are ascribed to beings of the world, but rather whether these properties are deemed fixed or malleable. On his view, we tend to conceptualize many inner properties as being fixed while indigenous Amazonian people conceptualize them as being highly malleable. In ethnobiology, essentialist thinking has often been associated with taxonomic thinking. Several authors have maintained that ethnobiological knowledge is ALIUS Research Group - Tribute to Martin Fortier ALIUS Bulletin n4 (2020) aliusresearch.org/bulletin 4 universally organized in taxonomic trees. In line with the Bayesian proposal, Martin explored the hypothesis that a given domain can be organized around various non-essentialist structures. His preliminary findings suggested that in some Amazonian cultures, abstract knowledge regarding biological beings is organized around complex and dynamical networks rather than taxonomic trees. Martin also investigated color categorization in Amazonian cultures. Some authors claim that color categories are organized according to universal principles in every culture around the world. Martin's fieldwork with the Shipibo-Konibo brought to light five basic color terms that work consistently with the color categories described in the World Color Survey. However, it also revealed the existence of a myriad of ad hoc color terms. For example, to refer to red, the Shipibo-Konibo often use the word "jimi" ("blood") even though they have a basic color term specifically referring to red ("joshin"). The ad hoc color terms used by the Shipibo-Konibo are polysemous (i.e., they can be used to refer to different colors if the original anchor object features different colors) and if speakers are asked about their meaning, they will first mention non-color meanings (e.g., "jimi" will be said to mean "blood" before "red"). These two properties contrast with lexicalized basic color terms used in English that originate in real objects of the world (e.g., the word "orange" for the orange color). In sum, Martin's work on color categorization among the Shipibo-Konibo questions the universality of classic color categorization, and suggests that some indigenous cultures of Amazonia have developed categorization systems that strikingly contrast with ours in that they include highly context-sensitive and flexible concepts. Many psychologists have claimed that persons are universally conceived in a dualist fashion. Martin investigated how the Shipibo-Konibo ascribe abilities ("thinking", "eating", "crying", "reasoning", etc.) to different beings (animals, plants, natural inanimate objects, artefacts); and also, how they ascribe "ontological components" ("shinan" = "thinking"; "kaya" = "soul"; "yora" = "body", etc.) to the same beings. His preliminary findings contrast with the view that dualism is universal. Instead, it seems that properties of beings are ascribed in a very dynamic and context-sensitive way. ALIUS Research Group - Tribute to Martin Fortier ALIUS Bulletin n4 (2020) aliusresearch.org/bulletin 5 This work also speaks to recent anthropological discussions around the concept of animism. As mentioned above, some anthropologists have it that a key feature of Amazonian cultures is that they over-ascribe mental properties to animals, plants and even non-living objects: in other words, there are animists. Martin's fieldwork on components of beings paints a different picture, as he did not find the kind of hyper-mentalization described by proponents of the animistic model. The Amazonian world is replete with "supernatural" entities, some of which are spirits and others described as masters or guardians of the animals, the river, the land, etc. (some supernatural entities are purported to be both spirits and masters). Martin studied the properties Shipibo-Konibo ascribe respectively to the master (ibo) of a plant (the ayahuasca vine), of a river (the Ucayali river), and of an animal (a white-lipped peccary). Interestingly, while the master of the plant is conceived as an immaterial spirit and as being the same for all the plants, the masters of the rivers and of the animals are taken to be other animals and each herd or each river possesses its own animal master. In the zoological and ecosystemic domains, the master seems to be defined by very specific ecological relationships. This work further questions some of the assumptions of approaches from both anthropology and the cognitive science of religion that over-emphasize the role of hyper- mentalization in folk intuitions about "supernatural" entities. An influential theory developed by several authors working on the cognitive science of religion claims that the main mechanisms of supernatural thinking lie in the violation of modular and universal domain-specific knowledge. In his work, Martin pointed out some weaknesses of this theory and developed an alternative proposal. In his view, supernatural thinking is triggered by violations of randomness that cannot be explained by our ordinary causal schemas. This theory is largely inspired from previous work on randomness, on probabilistic models of agency and on causal reasoning, and it contends that supernatural thinking is not a matter of modularity or domain-specific knowledge. This proposal challenges many influential approaches to religious thinking developed within the cognitive science of religion. ALIUS Research Group - Tribute to Martin Fortier ALIUS Bulletin n4 (2020) aliusresearch.org/bulletin 6 The second broad line of research that Martin was pursuing concerns about the relationship between hallucinogens, culture, and neurobiology. Psychotropic substances are pharmacological compounds that have the power of altering cognition and/or consciousness. Hallucinogens are a subclass of psychotropic compounds: they specifically alter perception by generating (visual, auditory, haptic, etc.) contents which are not veridical. Importantly, many psychostimulants (e.g., coca, khat) or sedative substances (e.g., opiates, kava) are not hallucinogenic. Martin's research was only concerned with hallucinogenic substances rather than psychotropic substances in general. After extensively reviewing the literature on hallucinogenic compounds and their effects, Martin identified nine classes of hallucinogens: serotonergic hallucinogens (e.g., LSD, psilocybin); antimuscarinic hallucinogens (e.g., scopolamine, atropine); antihistaminergic hallucinogens (e.g., brompheniramine, diphenhydramine); -opioid hallucinogens (e.g., pentazocine, bremazocine); antiglutamatergic hallucinogens (e.g., ketamine, phencyclidine); cannabinoid hallucinogens (e.g., THC, nabinole); nicotinic hallucinogens (e.g., nicotine, nornicotine); anticholinesterasic hallucinogens (e.g., harmine, harmaline); and GABAergic hallucinogens (e.g., muscimol, THIP). In addition to these nine classes, he subsumed compounds whose effects tap into several neurotransmission systems, such as ibogaine and noribogaine, into a tenth class. While most of the research investigating hallucinogens focuses on serotonergic compounds, Martin argued that critical insights on the mechanisms of drug-induced hallucinations can be gained by comparing different classes of hallucinogens. To that end, he undertook a general comparative study of the nine classes of hallucinogens, with a particular focus on the comparison between serotonergics and antimuscarinics, because these two are the most well documented. By comparing the underlying pharmacodynamics and resulting hallucinogenic effects of serotonergics and antimuscarinics, he suggested that one-to-one mappings can be established between the neuropharmacological and the phenomenological levels. ALIUS Research Group - Tribute to Martin Fortier ALIUS Bulletin n4 (2020) aliusresearch.org/bulletin 7 After showing how distinct neuropharmacological mechanisms correlate with distinct phenomenological contents, Martin proposed several neurocognitive models building the bridge between the cellular and the experiential levels. For example, a key feature of antimuscarinic hallucinations is that the experience is perfectly realistic and embedded in the non-hallucinatory environment. By contrast, serotoninergic hallucinations are mostly non-realistic and extraordinary in character and they present themselves as belonging to another "plane of reality". This phenomenological contrast can be readily explained by the fact that serotonergics importantly alter the early visual cortex as well as areas underlying object construction, while antimuscarinics mainly interfere in higher areas underlying scene construction. Martin believed that the investigation of the neuropharmacological and cognitive neuroscience of hallucinogens could also provide crucial insights for our understanding of hallucination and perception. In particular, he sought to debunk three "myths" endorsed by classical theories of hallucination and perception. The first myth is that genuine perception can be defined as a relation between the agent and the world, whereas hallucinations are characterized by the lack of such relation. Martin's work suggests that hallucinatory states are always influenced by inputs from the world. The real difference between perception and hallucination does not lie in whether the activity of central nervous system is related to the world (influenced by the world) or not, but in how stimuli from the world are being processed. A hallucinatory state may be induced exactly by the same input (i.e., by the same relation) as the non- hallucinatory state, but what makes it hallucinatory is the way it processes that input - i.e., what kind of output it generates given a certain input. The second myth is that there is such a thing as a perfect hallucination - i.e., a hallucinatory content that perfectly mimics a non-hallucinatory content. Martin argued that none of the different kinds of hallucination induced by his nine classes of hallucinogens ever generates perfect hallucinations. He believed this to be equally true of other types of hallucinations caused by ALIUS Research Group - Tribute to Martin Fortier ALIUS Bulletin n4 (2020) aliusresearch.org/bulletin 8 schizophrenia, neurodegenerative diseases, dreaming, etc. While hallucinations never perfectly mimic non-hallucinatory states, they may subjectively seem to do so. How is it, then, that subjects can erroneously experience a hallucinatory state to be identical to a non-hallucinatory state while it is not? In Martin's view, this happens every time the subject lacks metacognitive insights. This metacognitive faculty, he argued, is closely related to the activation of the dorsolateral prefrontal cortex. The third myth is that hallucinations are anomic, that is, that they are random (i.e., they can result in any kind of content and they are not governed by specific and predictable laws). This myth has been traditionally endorsed by authors such as Descartes and Malebranche. Martin's research suggests that each hallucinogen generates a specific kind of hallucinatory content and that this specificity can be elucidated by studying the cellular and neurocognitive mechanisms underpinning the hallucinatory effects. Thus, an antimuscarinic hallucinogen - unlike a serotonergic hallucinogen - will never induce extraordinary hallucinatory objects (such as a chimera). Thus, Martin argued that hallucinations are normatively constrained: they are governed by specific laws and these laws can be discovered by the neuroscience of hallucinations. On the basis of his investigation of real hallucinations, Martin rejected both "disjunctivism" and "common factor theories" (i.e., most brands of intentionalism, sense-data, etc.) in the philosophy of perception. He rejected disjunctivism because it claims that the key difference between the perceptual state and the hallucinatory state is to be understood in terms of relation to the world (present in one case and absent in the other). This amounts to endorsing the first myth. Conversely, he rejected "common factor theories" because they endorse the second myth of "perfect hallucinations". For that matter, he also rejected predictive coding accounts of hallucinations, which he believed to be also committed to the truth of the second myth (see the slogan "perception is a controlled hallucination"). Martin was also interested in the relationship between neuropharmacology and culture. He set out to answer the following question: when one entertains ALIUS Research Group - Tribute to Martin Fortier ALIUS Bulletin n4 (2020) aliusresearch.org/bulletin 9 a given hallucinogenic experience, what part of the experience is caused by neuropharmacology and what part is caused by culture? From a purely biological perspective, the whole experience might be boiled down to the properties of the pharmacological compound. Conversely, from a purely culturalist perspective, the whole experience might be said to be determined by one's mindset and by the setting in which the experience is taking place. Martin's contribution to this debate consists of highlighting two main points. First, he suggested that by comparing the role of neuropharmacology and culture in hallucinogenic experiences induced by distinct classes of compounds (serotonergics, antimuscarinics, -opioids, etc.), we can improve our understanding of the relative contribution made by each factor. Second, he argued that the neuropharmacology vs. culture divide is misleading because the two factors are in fact deeply intertwined. For example, he claimed that it is precisely for neuropharmacological reasons that experiences induced by serotonergic hallucinogens are so sensitive to the subject's mindset and the setting. The cultural shaping of such hallucinogenic experiences is constrained by neuropharmacology. Such cases require us to go beyond the usual neuropharmacology vs. culture divide and to develop more interactive and dynamic concepts. A related issue is the relationship between neuropharmacology and cultural evolution. A number of authors have made strong claims regarding the role played by the consumption of hallucinogens in cultural evolution. According to them, beliefs, artifacts and institutions are crucially shaped by hallucinogenic use. In other words, hallucinogenic use results in the emergence of very specific cultural traits. To assess such claims, Martin reviewed the key cultural traits these authors take to be specifically related to hallucinogenic use. In addition, he compared traditional cultures in which hallucinogenic use was present to others in which it was absent, and examined whether the latter lack the cultural traits purported to be found only in hallucinogenic-driven societies. This investigation revealed that the cultural traits under scrutiny are encountered in all cultures regardless of whether hallucinogenic use is present or not. Martin concluded that hallucinogens play only a limited role in cultural evolution; furthermore, ALIUS Research Group - Tribute to Martin Fortier ALIUS Bulletin n4 (2020) aliusresearch.org/bulletin 10 their role may concern formal properties of culture, rather than its content. To take the example of beliefs, Martin thought that hallucinogens shape not so much what is believed - i.e., the content of beliefs - as the kind of attitude people entertain towards those contents. In other words, regardless of what belief is being entertained (believing in gods, in spirits, in dwarfs, etc.), hallucinogens have very specific attitudinal effects. Van Leeuwen has proposed to distinguish religious credence from factual belief, and he (with others) argued that the attitude defined as factual belief is never found among religious believers. By contrast, Martin argued that hallucinogenic use precisely enables people to switch from religious credence to factual belief and make people believe in a "factual belief mode" what they previously believed in a "religious credence mode". Hence Martin's more general claim that the role hallucinogens play in cultural evolution is more a matter of attitudinal change than a matter of content change. In order to further investigate the interactions between hallucinogenic use and culture, Martin started developing an extensive database of hallucinogenic use across time and space, including hundreds of references from academic articles and books, explorers' diaries, and historical chronicles. Each of these references specifies which hallucinogen was used, where it was used, who was allowed to use it, what it was being used for (e.g., divinatory, political, or healing purposes), etc. This HUTHAC (Hallucinogenic Use Through History and Across Cultures) database was meant to be the first academic source exhaustively documenting hallucinogenic use through history and across cultures. The criteria used to insert data within the HUTHAC were particularly stringent. First, data included in the database concern only hallucinogenic use and not psychotropic use in general. Second, when building the HUTHAC, Martin made a distinction between use of hallucinogens at hallucinogenic doses and use of hallucinogens at sub-hallucinogenic doses (i.e., at doses too low to induce any hallucination). Sub-hallucinogenic uses of hallucinogens are not included in the HUTHAC. Since the 1950s, outlandish philological speculations have been made about the use of hallucinogens at different epochs and in different places in the world. The HUTHAC does not include the doubtful cases of hallucinogenic use advanced by these authors. It only ALIUS Research Group - Tribute to Martin Fortier ALIUS Bulletin n4 (2020) aliusresearch.org/bulletin 11 includes uses of hallucinogens for which there is a good amount of evidence. This evidence may be first-person observation (by an ethnographer or an explorer), third-person report (provided that the report can be considered reliable), and archeological evidence. The main preliminary finding of Martin's research on the HUTHAC database is that hallucinogenic use is fairly rare through history and across cultures. By and large, in the last hundred years, hallucinogenic use has mainly flourished in the Americas and been absent elsewhere (with only a few exceptions). Even in the Americas, hallucinogenic use has been persistently absent from some areas - notably North America (with the exception of Southern California and the Southwest), the Southern Cone, and Eastern Brazil. Another finding is that hallucinogenic use has gradually increased through the centuries (for example, in South America, ayahuasca and mimosa use had developed only in the 18th century). These findings challenge the view that hallucinogenic use is very ancient and belongs to an ancestral shamanic tradition that would have migrated from Siberia to the New World. More broadly, they also challenge the view that most archaic cultures used hallucinogens. Instead, they support the hypothesis defended by John Cooper and Joseph Wilbert, according to which hallucinogenic use is relatively recent and related to horticultural and agricultural practices. Some neuroscientists argue that our brain is evolutionarily wired, and as a result, it is everywhere the same and generates everywhere similar behavioral outputs. Other neuroscientists argue that the brain is wonderfully plastic and that its wiring can be extensively modified through experience. Martin believed that the study of hallucinogenic use provides us with a view of the brain that radically differs from these two stories. Regardless of whether our brain is malleable or not, what he found so fascinating about hallucinogens is that they can deeply modify behavior and experience, without modifying anything in the underlying hardware. When someone takes LSD, for example, the wiring of the brain is exactly the same before and during the hallucinogenic experience (one can certainly observe dramatic functional changes in brain connectivity, but not any structural change). And yet, under LSD, consciousness will be so deeply altered, sounds will influence the ALIUS Research Group - Tribute to Martin Fortier ALIUS Bulletin n4 (2020) aliusresearch.org/bulletin 12 content of vision, spatial dimensions will collapse, outlandish and folkloric characters will appear, new metaphysical beliefs will be endorsed, etc. Martin invented the concept of Homo neurochemicus to highlight that most of the flexibility of human behavior stems from neurochemical changes. He believed that proponents of evolutionary psychology and of fixist models of the wiring of the brain are mistaken in that they overlook the power of neurochemistry. All the behaviors and thoughts that they take to be universal and fixed can in fact be altered by introducing the right chemical compound in the central nervous system. According to him, proponents of the plasticity of the brain are also mistaken because they think that diversity of behavior and experience can only stem from a change in the wiring of the brain; in fact, functional changes caused by neurochemical means are much more dramatic than those induced by the long-term plasticity of the brain. It is obvious from this brief summary of Martin's work that he was a profoundly original and prolific researcher, promised to a brilliant academic career that would defy traditional disciplinary boundaries. His premature passing is all the more tragic that he cared deeply, and above all else, about answering the questions raised in these various projects. Beyond his extensive research program, Martin was a lion-hearted colleague and loyal friend. His contribution to academia was also pedagogical. During his stay in Peru, he dedicated a lot of time and effort to mentor undergraduate students to study abroad. Some of them have now entered Master studies and Ph.D. programs in top-ranked universities worldwide. He was able to connect people across continents and ideas across disciplines. Martin's legacy includes the formation of innovative collaborations (such as ALIUS), which will forever bear his unique signature. He was able to connect people across continents and ideas across disciplines. Martin's legacy includes the formation of innovative collaborations (such as ALIUS), which will forever bear his unique signature. His passion for the scientific study of consciousness was inspiring, altering the trajectory of many researchers in ALIUS Research Group - Tribute to Martin Fortier ALIUS Bulletin n4 (2020) aliusresearch.org/bulletin 13 our network and beyond. Now Martin lives on through networks: of neurons, citations, and colleagues. Selected references Fortier, M., Wente, A. O., Fernandez Flecha, M., \& Gopnik, A. (submitted). Abstract causal knowledge in Amazonia: Shipibo-Konibo horticulturalists learn multiple causality more readily than U.S. and Peruvian undergraduates. Fortier-Davy, M., \& Milliere, R. (2020). The multi-dimensional approach to drug- induced states a commentary on Bayne and Carter's "dimensions of consciousness and the psychedelic state". Neuroscience of Consciousness, 2020(1), niaa004. https://doi.org/10.1093/nc/niaa004 Fortier, M. (2020). Enculturation without TTOM and Bayesianism without FEP: Another Bayesian theory of culture is needed. Behavioral and Brain Sciences, 43. https://doi.org/10.1017/S0140525X19002905 Fortier, M. (2019). Le faconnement neuropharmacologique de la culture : Anthropologie comparee des rituels a hallucinogenes serotoninergiques et anticholinergiques. Cahiers d'Anthropologie Sociale, 17, 118-137. https://doi.org/10.3917/cas.017.0132 Fortier, M. (2019). L'ontologie amazonienne: animisme ou non-essentialisme? https://doi.org/10.31235/osf.io/7ntgh Fortier, M., Kellier, D., Flecha, M. F., \& Frank, M. C. (2018). Ad-hoc pragmatic implicatures among Shipibo-Konibo children in the Peruvian Amazon. https://doi.org/10.31234/osf.io/x7ad9 Fortier, M. (2018). Believing and experiencing the supernatural: Four shortcomings of the Interactive Religious Experience Model. Religion, Brain \& Behavior, 35- 38. https://doi.org/10.1080/2153599X.2018.1453531 Fortier, M. (2018). Le sens de realite dans les experiences psychotropes : Etude comparee des hallucinogenes serotoninergiques et anticholinergiques. Ed. S. Baud. Histoires et usages des plantes psychotropes. Paris: Imago, 125-184. Fortier, M. (2018). Sense of reality, metacognition, and culture in schizophrenic and drug-induced hallucinations. In J. Proust \& M. Fortier (Eds.), Metacognitive diversity: An interdisciplinary approach. Oxford/New York: Oxford University Press, pp. 343-378 ALIUS Research Group - Tribute to Martin Fortier ALIUS Bulletin n4 (2020) aliusresearch.org/bulletin 14 Proust, J., \& Fortier, M. (2018). Metacognitive diversity across cultures. In J. Proust \& M. Fortier (Eds.), Metacognitive diversity: An interdisciplinary approach. Oxford/New York: Oxford University Press, p. 1 Proust, J., \& Fortier, M. (Eds.). (2018). Metacognitive diversity: An interdisciplinary approach. Oxford University Press. Fortier, M., \& Kim, S. (2017). From the impossible to the improbable: a probabilistic account of magical beliefs and practices across development and cultures. In The Science of Lay Theories (pp. 265-315). Springer, Cham. Fortier, M., \& Fortier, M. (2017). Pour une science de la diversite de la conscience: Vers une redefinition multidimensionnelle des etats modifies de conscience. Intellectica, 67, 27-62. Dumas, G., Fortier, M., \& Gonzalez, J. (2017). Les enjeux des etats modifies de la conscience et de la cognition: limites passees et emergence de nouveaux paradigmes. Intellectica, 67, 27-62.\par
  \endgroup
}%
\end{document}
